\section{Design Rationale and Application Scope}
\label{sec:4}

\subsection{Define an observable support need}

We propose to begin with one event class, such as an unsuccessful attempt to contribute task-relevant information. The final definition must specify what constitutes an attempt, what evidence remains unshared, the eligible response interval, and conditions under which waiting is appropriate. This operationalization follows the construct-to-measure concern raised in multimodal collaboration research and avoids labeling a physiological value itself as a social need ~\cite{p57,p56}.

\missing{[TO COMPLETE: select the primary event; provide the task-specific definition, positive and negative examples, allowable response interval, and exclusion cases.]}

We propose triangulating independent behavioral coding, task logs, and event-based participant recall. Shared-arousal and medical-dyad studies illustrate the value of aligning bodily and conversational evidence with activity, while studies of biosignal interpretation show why participants' accounts should be retained ~\cite{p18,huang,p21}. Disagreement between observation and recall will be represented explicitly rather than resolved by treating either source as infallible.

\subsection{Three settings with distinct success criteria}

General XR collaboration will be the primary validation setting. The proposed task gives members different evidence that must be integrated into a shared decision. Candidate outcomes are relevant information contributed, acknowledged, and used, together with task quality and unwanted interventions. Moderation and attention-guidance work motivates the action options, but does not justify equal speaking time as the universal objective ~\cite{p70,p34}.

Medical team training will be a subsequent extension focused on communication under role and scenario constraints. Candidate targets include a critical observation that was not acknowledged or an unresolved communication breakdown. TEAM-related research and VR nontechnical-skills systems provide domain assessment foundations ~\cite{p13,p66,p4}. The extension will compare immediate prompting with preserving an episode for debriefing where appropriate; debriefing research motivates reflection, while the learning-performance distinction motivates a later scenario without prompts ~\cite{p9,p26,p62}. Trainee physiology and simulated patient physiology will remain separate data sources with separate meanings.

Collaborative education will be a subsequent extension focused on peer explanation and checking shared understanding. Candidate actions invite an explanation, ask peers to compare evidence, or request an individual prediction before discussion. These proposals draw on socially shared regulation, pedagogical-agent support, and collaborative VR scaffolding ~\cite{p29,p27,p51}. Outcomes will include individual understanding and unassisted transfer, because a correct group answer does not by itself demonstrate each member's retained capability ~\cite{p10,p62}.

\missing{[TO COMPLETE: confirm the primary domain and task; identify the medical and education partners, available participants, scenarios, and which extension is in scope for this paper.]}

\subsection{Design requirements}

D1 Ground decisions in the task. Candidate support needs must be interpretable in relation to roles, information access, task events, and recent interaction, as motivated by context-sensitive collaboration and regulation research ~\cite{p57,p29,huang}.

D2 Preserve uncertainty and choice. The agent must be able to wait, ask, or accept refusal; expectations research and human-AI interaction guidelines motivate configurable initiative and user correction ~\cite{expect,p3}.

D3 Separate sensing from action. A good estimate does not guarantee a useful intervention. Prediction, timing, acceptance, and consequences will be evaluated separately, drawing on proactive timing and intervention studies ~\cite{sensing,speaker,calm}.

D4 Treat physiology as fallible evidence. Signal quality, inter-person differences, and within-person context shifts require quality checks and robust comparisons ~\cite{artifacts,hhiss,crop}.

D5 Test capability beyond assisted performance. Training and education extensions require new tasks without prompts and, where feasible, delayed assessment ~\cite{p62}.
